\documentclass[a4paper,10pt, titlepage]{scrreprt}
%basic page

\linespread{1.5}				%Zeilenabstand 1,5mm
\parindent0cm						%Absatz wird nicht reingerückt

\usepackage[utf8]{inputenc}
\usepackage{graphicx}				%Zum Einbinden von externen Grafiken (jpeg,pdf)
\usepackage{a4wide}					%Ausnutzen der ganzen Blattbreite
\usepackage{fancyhdr} 			%Package für Kopf- und Fusszeile
\usepackage[ngerman]{babel}	%Rechtschreibung
\usepackage{hyperref}				%Einfügen von Weblinks
\usepackage{amsmath}				%Paket für mathematische Formeln
\usepackage{amsfonts}				%zusätzliche mathematische Zeichen
\usepackage{amssymb}				%zusätzliche mathematische Symbole
\usepackage[labelfont=bf]{caption}  %"Abbildung" fett geschrieben
\usepackage{subfigure} 			%Bilder da plazieren wo sie auch im Latex-Code stehen
\usepackage{graphicx}
\usepackage{eso-pic,picture}
\usepackage[absolute]{textpos}
\usepackage{hyphenat}
\usepackage[onehalfspacing]{setspace}
\usepackage{multirow}
\usepackage{array}
\usepackage{caption}
\usepackage{url}
\usepackage{chngcntr}
\usepackage{tikz}
\counterwithin{figure}{section}
\usepackage[a4paper, left=3.5cm, right=2.5cm, top=2.0cm, bottom=2.0cm,footskip=1.5cm, includeheadfoot] {geometry}  % manuelles Seitenlayout
\usepackage{bibgerm}        % Einfügen des Literaturverzeichnisses
\usepackage{eurosym}        % Einfügen des € Zeichens
%\usepackage{lscape}        % Text einer Seite wird im quer dargestellt
%\usepackage[a4paper]{geometry}       % Querformat einer Seite
\usepackage{pdfpages}       % Einfügen von ganzen PDF-Seiten
\usepackage{array}

%================================================================================

%% definieren einer neuen Tabelle mit automatischen Umbruch
\newcolumntype{C}[1]{>{\centering\arraybackslash}m{#1}}%m-vertikal zentriert, b bottom, p top; allg: zur Definition einer festen Spaltenbreite
\newcolumntype{L}[1]{>{\raggedright\arraybackslash}p{#1}} % linksbündig mit Breitenangabe


%================================================================================


\begin{document}

%Titelseite
%--------------------------------------------------------------------------------
\begin{titlepage}
\thispagestyle{empty} 		%Keine Kopf- und Fusszeilen, keine Seitenzahl
\linespread{1.0}

\begin{figure}
\includegraphics[scale=0.3]{images/HTW-Logo}
\end{figure}

\hrule
\vspace{0.2 cm}
\textbf{Fakultät Elektrotechnik}
\vspace{0.1 cm}
\hrule
\vspace{0.2 cm}
Studiengang: Mechatronik/Fahrzeugmechatronik 

\begin{flushleft}
\vspace{8 cm}
\huge\textbf{Praktikumsbericht}\\ [0.5 cm]
\Large{Bericht zur Praktikum im Adensis GmbH}\\ [1.5 cm]
\Large{Dresden, DATUM}
\end{flushleft}
\pagebreak 

%new page

\thispagestyle{empty} 		%Keine Kopf- und Fusszeilen, keine Seitenzahl
\begin{figure}
\includegraphics[scale=0.3]{images/HTW-Logo}
\end{figure}

\hrule
\vspace{0.2 cm}
\textbf{Fakultät Elektrotechnik}
\vspace{0.1 cm}
\hrule
\vspace{0.2 cm}
Studiengang: Mechatronik/Fahrzeugmechatronik 

\begin{center}
\vspace{2 cm}
\huge\textbf{PRAKTIKUMSBERICHT}
\vspace{2 cm}
\end{center}

\begin{tabular}{lL{7cm}}
\Large\textbf{Thema:} & \Large{Praktikumsbericht}\\
\Large\textbf{Bearbeiter:} & \Large{Oskar Engler}\\
\Large\textbf{Matrikelnummer:} & \Large{34431}\\
\Large\textbf{Bearbeitungszeitraum:} & \Large{01.10.14 bis 28.02.15}\\
\Large\textbf{Ort, Datum der Abgabe:} & \Large{Dresden, DATUM}\\
\Large\textbf{Betreuer:} & \Large{Dipl. Ing. Lars Mademann ??}\\
\Large\textbf{Verantw. Hochschullehrer:} \hspace{1 cm} & \Large{Prof. Dr.-Ing. Ralf Boden}\\
\vspace{1 cm}\\
\large{Textseiten:} & \large{xx}\\
\large{Anlagen:} & \large{yy}\\
\large{Anhänge:} & \large{zz}\\
\end{tabular}

\end{titlepage}
\linespread{1.5}

%Aufgabenstellung
%---------------------------------------------------------------------------------
\fancyhf{}

\thispagestyle{fancy}
\fancyhead[L]{Oskar Engler}
%Kopfzeile mittig
\fancyhead[C]{Matr-Nr.: 34431}
\fancyhead[R]{HTW-Dresden}
%\renewcommand{\headrulewidth}{0.5pt}

\pagebreak 

%Hauptteil
%---------------------------------------------------------------------------------
%Kapitel sind für eine bessere Übersicht in eigenen Dateien und werden hier aufgerufen:

\chapter*{Sperrvermerk}\thispagestyle{fancy}\markboth{Sperrvermerk}{1}
\addcontentsline{toc}{chapter}{Sperrvermerk}

Diese Praktikumsbericht enthält vertrauliche Informationen, die der Geheimhaltung unterliegen. Sie dürfen nur für die interne Verwendung und zur Kontrolle durch den verantwortlichen Hochschullehrer genutzt werden. Eine, auch nur teilweise, Veröffentlichung der Belegarbeit darf nur mit Zustimmung der BELECTRIC GmbH, Zweigstelle Dresden, Industriestraße 65, 01129 Dresden erfolgen.
\\ 
\\
\\
{\LARGE Dresden, DATUM}\\


%Kopfzeile
%-------------------------------------------------------------------------------- \fancyhf{} %alle Kopf- und Fußzeilenfelder bereinigen
\pagestyle{fancy} %eigener Seitenstil

%Kopfzeile links bzw. innen
\fancyhead[L]{Oskar Engler}
%Kopfzeile mittig
\fancyhead[C]{Matr-Nr.: 34431}
%\facnyhead[R]{\chaptermark}
\renewcommand{\headrulewidth}{0.5pt}				%obere Linie

%Fußzeile
%--------------------------------------------------------------------------------
\fancyfoot[R]{\thepage} 							%Seitennummer
\fancyfoot[L]{\nouppercase{\leftmark}}							%Kapitelbezeichnung Position
\renewcommand{\chaptermark}[1]{\markboth{\thechapter.~\chaptername: #1}{}}	%Kapitelbezeichnung 

\setcounter{page}{1}
\pagenumbering{Roman}
%Inhaltsverzeichnis plotten
%---------------------------------------------------------------------------------
\addtocontents{toc}{\protect\thispagestyle{fancy}}
\thispagestyle{fancy}
\tableofcontents
\thispagestyle{fancy}
% Abkürzungsverzeichnis
%---------------------------------------------------------------------------------

\chapter*{Abkürzungsverzeichnis}\thispagestyle{fancy}\markboth{Abkürzungsverzeichnis}{Abkürzungsverzeichnis}
\addcontentsline{toc}{chapter}{Abkürzungsverzeichnis}

\begin{tabbing}
\ \= \textbf{Abkürzung} \hspace{0.5cm} \= \textbf{Beschreibung} \kill
\\

\> BMS \>   Batterie-Management-System\\
\> EBU \>   Energy Buffer Unit\\
\> B\&R \>		Bernecker + Rainer Industrie-Elektronik\\
\> ST \>   Strukturierter Text\\
\> SQL \>   Structured Query Language\\
\> HTML \>  HyperText Markup Language \\
\> CSS \>   Cascading Style Sheets\\
\> PHP \>   Hypertext Preprocessor\\
\> CSV \>   Comma-separated values\\
\> JSON \>  JavaScript Object Notation\\
\> FTP \>   File Transfer Protocol\\



\end{tabbing}

% Symbolverzeichnis
%---------------------------------------------------------------------------------

\chapter*{Symbolverzeichnis}\thispagestyle{fancy}\markboth{Symbolverzeichnis}{Symbolverzeichnis}
\addcontentsline{toc}{chapter}{Symbolverzeichnis}

\section*{Formelzeichen}

\begin{tabbing}
\ \= \textbf{Symbol} \hspace{0.5cm} \= \textbf{Einheit} \hspace{1.5cm} \= \kill
\> $A$ \> $\mathrm{m^{2}}$ \> Fläche \\
\> $a$ \> $\mathrm{MPa}$ \> Kohäsionsdruck \\
\end{tabbing}

\section*{Griechische Symbole}

\begin{tabbing}
\ \= \textbf{Symbol} \hspace{0.5cm} \= \textbf{Einheit} \hspace{1.5cm} \= \kill
\> $\alpha$ \> $\mathrm{W\,/\,(m^{2} \cdot K)}$ \> Wärmeübergangskoeffizient\\
\> $\gamma$ \> - \> Aktivitätskoeffizient für die Flüssigphase \\
\end{tabbing}
% Abbildungsverzeichnis
%---------------------------------------------------------------------------------
\setcounter{lofdepth}{-1}
\listoffigures
\setcounter{lofdepth}{-1}

\end{document}